\documentclass[11pt]{article}
\usepackage[margin=1in]{geometry}
\usepackage[T1]{fontenc}
\usepackage{lmodern,microtype,hyperref,enumitem,amsmath,amssymb,xcolor}
\usepackage{listings}
\hypersetup{colorlinks=true,urlcolor=blue,citecolor=blue,linkcolor=blue}

\lstset{basicstyle=\ttfamily\small,columns=fullflexible,keepspaces=true,breaklines=true}

\title{Q09: How Does the Composition of Training Data Shape What a Model Knows and Can Do?}
\author{DLS 2026 --- Module (Paradise)}
\date{October 2026}

\begin{document}
\maketitle

\section*{Learning objectives}
After this module, a student should be able to (i) describe what web-scale pretraining corpora actually contain and how curation pipelines (extraction, filtering, deduplication) shape them, (ii) explain exact and MinHash-based fuzzy deduplication and why dedup improves training efficiency and reduces memorization, (iii) reason about quality filtering (heuristic, classifier-based, and perplexity-based) and its trade-offs, (iv) state what is known about repetition and reuse (the $\sim$4-epoch result), mixture weighting, curriculum ordering, and tokenizer--data interactions, and (v) articulate why data quality is now treated as a first-class lever alongside raw scale.

\section{Introduction: from ``more data'' to ``which data''}
Scaling-law analysis made dataset size a headline number: compute-optimal training scales parameters and tokens together, with Chinchilla trained on $1.4$T tokens rather than the $\sim$300B of its same-compute predecessor Gopher \cite{hoffmann2022chinchilla}. But two findings pushed the field past ``just add tokens.'' First, extrapolating demand against the finite stock of public human text suggests the stock could be fully utilized between roughly 2026 and 2032 \cite{villalobos2024will}. Second, controlled ablations show that \emph{which} tokens you use changes benchmark performance by amounts comparable to modest scale increases: filtering the 15T-token FineWeb corpus down to 1.3T tokens of classifier-scored educational content (FineWeb-Edu) raised MMLU from about 33\% to 37\% and ARC from about 46\% to 57\% at a fixed 350B-token training budget \cite{penedo2024fineweb}. Composition --- sources, filters, duplicates, mixture weights, and ordering --- is now a competitive discipline, not bookkeeping. This note treats each lever in turn, staying honest about which effects are established at what scale.

\section{What is in a web-scale corpus?}
Modern pretraining sets are dominated by Common Crawl snapshots, plus curated slices: code, Wikipedia, books, and papers. The Stack collects 3.1\,TB of permissively licensed source code across 30 languages, and is the reference code corpus for open models \cite{kocetkov2022stack}. Web crawl, by contrast, arrives as raw HTML of highly uneven quality; every major pipeline therefore performs text extraction, language identification, URL blocklists, heuristic quality filters, and deduplication before any training \cite{penedo2024fineweb}. Domain composition matters in itself: pretraining 1.5B-parameter models on different domain mixtures shows that including heterogeneous sources (web plus books) is broadly beneficial, while single-source curricula are worse, and filter effects are not predictable from surface domain statistics \cite{longpre2023pretrainer}. The upshot: a ``token'' is not a fungible unit. A token from a deduplicated Wikipedia article, a scraped boilerplate page, and a Python function carry very different learning signal per FLOP.

\section{Deduplication}
\paragraph{Exact and near-duplicate removal.} Web corpora contain massive redundancy: one 61-word English sentence appeared over 60{,}000 times in C4, and removing near-duplicates and long repeated substrings let language models emit memorized text roughly ten times less frequently while reaching the same accuracy in fewer steps; deduplicated C4 is a few percent smaller and correspondingly cheaper per epoch \cite{lee2022deduplicating}. Exact dedup (hashing documents or lines) is cheap but misses the common case of boilerplate templates and paraphrased copies.

\paragraph{MinHash and fuzzy dedup.} Near-duplicate detection needs a similarity notion robust to small edits. The standard tool, MinHash, works on \emph{shingles}: fixed-length token $n$-grams sampled from a document. If a document is a set $S$ of shingles, the Jaccard similarity of two documents is $J(A,B) = |A \cap B| / |A \cup B|$. MinHash uses $k$ random hash functions $h_i$ and stores, per document, the signature $\big(\min_{s \in S} h_i(s)\big)_{i=1..k}$: the probability that two documents agree in one signature coordinate equals their Jaccard similarity. Signatures are bucketed via locality-sensitive hashing (banding) so only candidate pairs are compared, which is what makes this subquadratic at trillion-token scale. FineWeb ablated dedup granularity extensively; a notable result is that \emph{global} MinHash dedup across all 96 snapshots removed $\sim$80\% of data (from 20T to 4T tokens) and trained \emph{worse} than per-snapshot dedup, which kept 20T tokens and matched RefinedWeb --- aggressive dedup can discard useful variety \cite{penedo2024fineweb}. Muennighoff et al.\ similarly report that dedup did not help downstream on their (already clean) C4-based setup, while noting it may still reduce memorization \cite{muennighoff2023scaling}.

\paragraph{Why dedup helps when it helps.} Three mechanisms: (1) \emph{efficiency} --- repeated documents waste gradient steps on already-seen content; (2) \emph{memorization and privacy} --- duplication is the strongest known driver of verbatim recall, so removing duplicates reduces regurgitation and extraction risk \cite{lee2022deduplicating}; (3) \emph{evaluation validity} --- train/test overlap inflates benchmark scores. The caveat above is real: these benefits concentrate on noisy corpora, and overly aggressive dedup is harmful \cite{penedo2024fineweb,muennighoff2023scaling}.

\section{Quality filtering}
Three families of filters, in historical order:
\begin{enumerate}[noitemsep]
  \item \textbf{Heuristics:} C4-style rules (document length, ratio of punctuation/lines starting with bullets, presence of stopwords). Cheap and blunt; FineWeb's ablations found that keeping a small set of effective C4/MassiveText-style filters, tuned rather than piled on, was a good default \cite{penedo2024fineweb}.
  \item \textbf{Perplexity filtering:} CCNet scores each document by the perplexity of an LM trained on a target domain (e.g., Wikipedia) and keeps low-perplexity text \cite{wenzek2020ccnet}. This is \emph{domain targeting} disguised as quality: it selects for Wikipedia-likeness. Muennighoff et al.\ found perplexity filtering effective on noisy data but not on clean data \cite{muennighoff2023scaling}.
  \item \textbf{Classifier filtering:} FineWeb-Edu labels $\sim$460K sampled pages with Llama-3-70B-Instruct for educational value (0--5), trains a regressor on embeddings, and keeps pages above a threshold \cite{penedo2024fineweb}. This ``LLM-as-annotator'' approach scales curation judgments, but inherits the judge model's biases: it optimizes \emph{grade-school-educational} text, and the resulting 8$\times$ reduction in tokens is only a win because quality-per-token rose faster than count fell.
\end{enumerate}
Filtering has a known dark side: aggressive quality or toxicity filtering trades benchmark performance against representation and safety properties, and no one-size-fits-all threshold exists; effects of filtering are not predictable from text statistics and must be ablated empirically \cite{longpre2023pretrainer}. The textbook-quality extreme --- small, heavily curated, partly synthetic corpora --- produced strong small models for code \cite{gunasekar2023textbooks}, demonstrating that quality can substitute for scale up to a point, with open questions about what such narrow data fails to teach.

\section{Mixture weighting, repetition, and annealing}
\paragraph{Domain mixtures.} Training sets upsample high-value domains (Wikipedia, code, papers) far above their natural web frequency. DoReMi shows domain weights can be \emph{optimized}: a small proxy model trained with group-DRO over domains yields weights that, applied to a 30$\times$ larger model, improved few-shot accuracy by 6.5 points over The Pile's default mixture and reached baseline accuracy in 2.6$\times$ fewer steps \cite{xie2023doremi}. Mixture choices interact with the target: code in the mix has repeatedly been observed to benefit natural-language reasoning; Muennighoff et al.\ quantify one version of this, finding that filling up to half the data budget with Python code (from The Stack) gave roughly a $2\times$ increase in \emph{effective} tokens even when evaluating only natural-language tasks \cite{muennighoff2023scaling,kocetkov2022stack}. The mechanism is hypothesized to be shared structure (long-range state tracking, compositional dependencies) rather than surface syntax, though we flag that code$\to$math transfer specifically is reported anecdotally across models rather than established by a single controlled study.

\paragraph{Repetition: the $\sim$4-epoch rule.} With 400+ runs up to 9B parameters, Muennighoff et al.\ find that with constrained data and fixed compute, up to about \textbf{4 epochs} of repeated data changes loss negligibly relative to unique data; beyond that, returns decay toward zero \cite{muennighoff2023scaling}. Their fitted law predicts that in the data-constrained regime extra compute should go disproportionately to epochs rather than parameters (smaller model, more passes) --- unlike Chinchilla's equal scaling \cite{hoffmann2022chinchilla}. Caveats: their results are for repetition of a \emph{whole} corpus from the start of training; up-sampling a single domain is a different intervention, and the resulting rule is a practical heuristic, not a constant of nature.

\paragraph{Annealing.} A common recipe schedules the mixture over training: a large, diverse mix early, then a shorter ``annealing'' phase near the end (with decaying learning rate) on upsampled high-quality data. MiniCPM's two-stage recipe raised the mathematics and code share to $\sim$70\% during the annealing phase of pretraining \cite{hu2024minicpm}. Reported across many models but rarely ablated in isolation at scale, annealing should be read as a converging practice with plausibly large effects and thin causal literature.

\section{Curriculum and ordering}
Beyond mixture schedules, the \emph{order} of data matters at the margins. Data age is one documented axis: a temporal gap between pretraining and evaluation data degrades performance, and the gap is not repaired by finetuning \cite{longpre2023pretrainer}. Within a fixed mixture, uniform random shuffling is the default because it keeps the batch distribution stationary; easy-to-hard curricula can beat it for noisy data (RHO-LOSS selects learnable, relevant, not-yet-learned points, accelerating training on web-scraped data \cite{mindermann2022prioritized}), but results are setup-dependent and rarely survive the web-scale setting unchanged. The honest summary: ordering effects are real, modest, and mostly exploited via coarse mixture schedules (annealing) rather than fine-grained curricula.

\section{Tokenizer and vocabulary interactions}
The data mixture also flows into the \emph{tokenizer}, which is trained on a sample of the corpus and fixes the tokenization for all of training. Consequences:
\begin{itemize}[noitemsep]
  \item \textbf{Fertility differences across domains:} code and non-English text split into more subword units per word than English prose, so a fixed token budget holds less code content; ``tokens'' as a unit of accounting are domain-relative. BPE-style subword tokenization \cite{sennrich2016neural} trades vocabulary size against sequence length, and its merge statistics reflect the tokenizer-training mixture.
  \item \textbf{Mixture accounting:} domain weights specified in tokens shift when measured in raw characters or bytes, since fertility differs by domain; a tokenizer trained on web text will fragment the scientific vocabulary of a papers-heavy mixture, spending capacity on rare subwords.
  \item \textbf{Vocabulary coverage:} a tokenizer trained without code (or without a language) makes code (or that language) artificially expensive, biasing the effective mixture away from what was declared.
\end{itemize}
A clean statement of the interaction is hard to find in the literature; the mechanism is uncontroversial but the magnitude of the effect relative to mixture weights is, to our knowledge, not separately quantified at scale --- a good open question.

\section{Why quality now rivals scale}
Putting the pieces together: (1) the marginal value of fresh web tokens falls as the easy stock is consumed \cite{villalobos2024will}; (2) repeated and code-augmented tokens remain useful, with quantified diminishing returns \cite{muennighoff2023scaling}; (3) classifier-based filtering shifts the quality-per-token distribution enough to change benchmark rankings at fixed compute \cite{penedo2024fineweb}; (4) optimized mixtures buy training-step efficiency comparable to nontrivial compute increases \cite{xie2023doremi}. None of these licenses the claim that ``data quality beats scale'' unconditionally --- Chinchilla's compute-optimal frontier still binds \cite{hoffmann2022chinchilla} --- but at fixed compute, composition is among the cheapest levers available, and it is the only lever whose gains compound with every downstream epoch.

\paragraph{What we did not establish.} This note deliberately hedges: code$\to$math transfer mechanisms, annealing ablations at scale, and tokenizer--mixture interaction magnitudes lack the controlled, multi-scale evidence base of scaling laws or the 4-epoch result. Treat those as strong practitioner priors, not settled science.

\bibliographystyle{plain}
\bibliography{bibliography}

\end{document}
